\section*{Chapter \ref{ch:g11:inequality} --- Inequality}
% figdata: figdata/grade-11/11-inequality/out/ (eu_gini, lorenz_gini, bren_fifths, bren_key); Brenland numbers invented.

\begin{solution}{exo:g11:inequality:1}
$(0, 0)$, $(20, 5)$, $(40, 15)$, $(60, 30)$, $(80, 55)$, $(100, 100)$, in per cent.
\end{solution}

\begin{solution}{exo:g11:inequality:2}
$(0, 0)$, $(20, 8)$, $(40, 20)$, $(60, 38)$, $(80, 62)$, $(100, 100)$.
\end{solution}

\begin{solution}{exo:g11:inequality:3}
Bulgaria, 37.7 after transfers; France has the largest fall, $53.7 - 30.4 = 23.3$ points. The EU: 47.8 before and 29.2 after.
\end{solution}

\begin{solution}{exo:g11:inequality:4}
$60\,000/15\,000 = 4$.
\end{solution}

\begin{solution}{exo:g11:inequality:5}
$20/100 = 20\,\%$; $50/100 = 50\,\%$.
\end{solution}

\begin{solution}{exo:g11:inequality:6}
$\sum (L_{k-1} + L_k) = 0.05 + 0.20 + 0.45 + 0.85 + 1.55 = 3.10$, so $G = 1 - 0.2 \times 3.10 = 0.38$.
\end{solution}

\begin{solution}{exo:g11:inequality:7}
$0.08 + 0.28 + 0.58 + 1.00 + 1.62 = 3.56$, so $G = 1 - 0.712 = 0.288 \approx 0.29$. The second country is more equal.
\end{solution}

\begin{solution}{exo:g11:inequality:8}
$8000 - 0 + 6000 = 14\,000$; $30\,000 - 3000 + 1000 = 28\,000$; $90\,000 - 21\,000 = 69\,000$. Richest to poorest: $90\,000/8000 \approx 11.3$ before, $69\,000/14\,000 \approx 4.9$
after.
\end{solution}

\begin{solution}{exo:g11:inequality:9}
Income tax: average rates 5, 10 and 15\,\%, rising with income: progressive. Purchase tax: 9, 6 and 4\,\%, falling: regressive.
\end{solution}

\begin{solution}{exo:g11:inequality:10}
Bottom row, of 100: $50/100 = 50\,\%$ bottom, 30\,\% middle, 20\,\% top. Top column, of $20 + 25 + 55 = 100$: 20\,\% have bottom-third parents, 25\,\%
middle, 55\,\% top. Each row and each column holds 100 families, a third of 300.
\end{solution}

\begin{solution}{exo:g11:inequality:11}
Total 20\,000; shares 5, 10, 15, 20, 50; cumulative 0.05, 0.15, 0.30, 0.50, 1; $\sum(L_{k-1} + L_k) = 3.00$, $G = 0.40$. After: total 25\,000, shares 8, 12, 16, 20, 44;
cumulative 0.08, 0.20, 0.36, 0.56, 1; sum 3.40, $G = 0.32$. The same sum added to every income lowers the Gini.
\end{solution}

\begin{solution}{exo:g11:inequality:12}
A rising premium while the supply of graduates grew faster than any plausible rise in demand, or a falling premium while supply stalled. The halt after 2004 is not such
evidence: it came when the supply of new graduates rose fast, as the model predicts. The rise of the top 1\,\% share is not evidence against the model either, but lies
outside it: the model is about two broad groups, and says nothing about the very top.
\end{solution}

\begin{solution}{exo:g11:inequality:13}
Before: total 200, shares 5, 10, 15, 20, 50, $G = 0.40$ (as in \cref{exo:g11:inequality:11}). After: total 200, shares 4.5, 12.5, 17.5, 20.5, 45; cumulative 0.045, 0.17,
0.345, 0.55, 1; sum $0.045 + 0.215 + 0.515 + 0.895 + 1.55 = 3.22$, $G = 1 - 0.644 = 0.356$. The Gini falls from 0.40 to 0.36, while the poorest household goes from 10 to 9.
\end{solution}

\begin{solution}{exo:g11:inequality:14}
Total income: B (36\,000) above A (20\,000). Per person: A 20\,000, B 9000. Square root: A 20\,000, B $36\,000/2 = 18\,000$. Modified OECD: A 20\,000, B $36\,000/2.1 \approx 17\,140$.
B is richer on the total, A on every scale that allows for size.
\end{solution}

\begin{solution}{exo:g11:inequality:15}
People save in good years and borrow or draw on savings in bad ones, so their spending varies less than their yearly income; students and retired people spend more than
they earn. The South African survey measures consumption, the others income, so its curve would likely lie closer to the diagonal than one of income would.
\end{solution}

\begin{solution}{pb:g11:inequality:1}
\textbf{1.} $36\,000/\sqrt{4} = 18\,000$: on this scale it is as well off as the single person.

\textbf{2.} People living together share housing, heating and much else, so four people do not need four times the income of one to live as well.

\textbf{3.} $1 + 0.5 + 0.3 + 0.3 = 2.1$; $36\,000/2.1 \approx 17\,140$ dalens.

\textbf{4.} $(0, 0)$, $(20, 6.9)$, $(40, 19.7)$, $(60, 37.2)$, $(80, 60.8)$, $(100, 100)$.

\textbf{5.} $B = 0.1 \times (0.069 + 0.266 + 0.569 + 0.980 + 1.608) = 0.1 \times 3.492 = 0.3492$.

\textbf{6.} $G = 1 - 2 \times 0.3492 = 0.302$, about 0.30.

\textbf{7.} $39.2/6.9 \approx 5.7$ times.

\textbf{8.} $(0, 0)$, $(20, 10.0)$, $(40, 23.9)$, $(60, 41.6)$, $(80, 64.5)$, $(100, 100)$.

\textbf{9.} $\sum (L_{k-1} + L_k) = 0.100 + 0.339 + 0.655 + 1.061 + 1.645 = 3.800$, so $G = 1 - 0.760 = 0.240$.

\textbf{10.} $0.302 - 0.240 = 0.062$: about 6 points on a scale of 0 to 100.

\textbf{11.} The first three fifths gain (by 3.1, 1.1 and 0.2 points); the fourth loses 0.7 and the top fifth 3.7.

\textbf{12.} Brenland's fall of about 6 points is smaller than every fall in the chart (13.8 to 23.3). Eurostat's \emph{before} removes pensions and other social transfers
but leaves income taxes deducted, while Brenland's is before taxes and transfers; and Eurostat computes the Gini from each household, not from fifths.

\textbf{13.} \omterm{def:g11:inequality:market-income}{Market income}: $33\,900/7800 \approx 4.3$. Disposable: $28\,600/9300 \approx 3.1$.

\textbf{14.} From 23.6 to 20.8\,\%: 2.8 points.

\textbf{15.} Gini: $0.062/0.302 \approx 21\,\%$. \omterm{def:g11:inequality:interdecile}{Interdecile ratio}: $(4.3 - 3.1)/4.3 \approx 28\,\%$. Top tenth: $2.8/23.6 \approx 12\,\%$. The \omterm{def:g11:inequality:interdecile}{interdecile ratio} falls most: taxes and
transfers do most at the bottom of the distribution.

\textbf{16.} Within each fifth incomes differ, and the straight lines between the five points hide that inequality; the true \omterm{def:g11:inequality:lorenz}{Lorenz curve} lies below them.

\textbf{17.} For example: services provided free (schooling, health care), and taxes on spending; it also says nothing about \omterm{def:g11:inequality:wealth}{wealth}.

\textbf{18.} It falls: each income becomes $(1 - t)y$ plus the same payment for all, which raises the shares of the poorer households and lowers those of the richer.

\textbf{19.} Yes: in \cref{exo:g11:inequality:13} the Gini falls from 0.40 to 0.36 while the poorest household loses.

\textbf{20.} The \omterm{def:g11:inequality:gini}{Gini coefficient} is 0.30 for \omterm{def:g11:inequality:market-income}{market income} and 0.24 for disposable income: taxes and transfers lower it by about 6 points.
\end{solution}
